\section{A chromatic high-minimum-degree condition}\label{sec:highdegree}

We use the following published degree-constrained subgraph theorem
of Addario-Berry, Dalal and Reed~\cite{ADR08}, Theorem~5 in the
author version. For a simple host graph with degrees $d(v)$, integers
$a_v,b_v$ satisfying
\begin{equation}\label{eq:ADR}
 0\le a_v\le\lfloor d(v)/2\rfloor\le b_v<d(v),\quad
 b_v\le(d(v)+a_v)/2+1,\quad b_v\le2a_v+3
\end{equation}
admit a spanning subgraph $F$ with
$d_F(v)\in\{a_v,a_v+1,b_v,b_v+1\}$. We quote this theorem and
prove the residue and boundary consequences below.

\begin{theorem}[D1]\label{thm:chromatic}
If a simple graph has a proper $c$-vertex-colouring, $c\ge2$, and
\[
 \delta(G)\ge T(c):=8c(2c+1)-7,
\]
then it admits doubled weights $x_e\in\{0,1\}$ for which both
colour degrees distinguish the ends of every original edge. The same
internal conclusion holds with arbitrary prescribed physical external
counts at every vertex.
\end{theorem}
\begin{proof}
Put $p=2c,q=2c+1,M=pq$, and let $j(v)\in\{0,\ldots,c-1\}$
be the proper colours. With external counts $A_v,B_v$ and internal
degree $d=d_G(v)$ choose a residue $\phi_v$ by
\[
 \phi_v\equiv2j(v)-A_v\pmod p,\qquad
 \phi_v\equiv B_v+2d-2j(v)\pmod q.
\]
The moduli are coprime. Set $t=\lfloor d/2\rfloor$,
$r=(t-\phi_v)\bmod M$, $a=t-r$, and $b=a$ if $r=0$, otherwise
$b=a+M$. Both $a,b$ have residue $\phi_v$ and bracket $t$.
The assumption is $d\ge4M-7$, so
$\lfloor d/2\rfloor\ge2M-4$ and $\lceil d/2\rceil\ge2M-3$.
For $r>0$,
\[
 a\ge M-3,\quad b=a+M\le2a+3,\quad
 2b\le d+a+2\ \Longleftrightarrow\
 \lceil d/2\rceil\ge2M-r-2,
\]
and the last inequality holds since $r\ge1$. Also
$b\le t+M-1\le d-1$. For $r=0$, all conditions are immediate.
Thus~\eqref{eq:ADR} holds in the original simple graph.

Give $F$ edges weight 1 and all others zero. For some
$\epsilon_v\in\{0,1\}$ the final red and blue residues are
\[
 R(v)\bmod p\in P_{j(v)}:=\{2j(v),2j(v)+1\},\qquad
 B(v)\bmod q\in Q_{j(v)}:=\{2j(v),2j(v)-1\}.
\]
The $P_j$ are pairwise disjoint modulo $2c$. The $Q_j$ are also
pairwise disjoint modulo $2c+1$: $Q_0=\{0,2c\}$, and the others
are $\{1,2\},\{3,4\},\ldots,\{2c-3,2c-2\}$. Different
proper colours thus give unequal actual degrees in both colours.
The choice of external offsets changes no degree inequality, proving
the interface version too.
\end{proof}

For $c=2,3,4,5$ the thresholds are $73,161,281,433$. They are
sufficient bounds, with no optimality claim. The theorem has been
applied to a simple host only: on two vertices with 100 parallel
edges, the parameter pairs $(25,50)$ and $(30,55)$ satisfy the
numerical inequalities~\eqref{eq:ADR}, but the allowed factor-degree
sets at the two ends are disjoint, so no such multigraph factor
exists. Substituting the doubled host into the quoted theorem is
therefore not justified by these numerical conditions.

\subsection{One internal boundary vertex}

Suppose a connected induced region $H$ has all its outside edges
at one vertex $w$. Fix every edge outside the interior of $H$ and
assume that outside--outside edges are legal using their final counts.
If $h$ boundary edges include $m$ mixed ones, a pure boundary edge
forbids at most one residue label and a mixed edge at most two,
since the $P_j$ and $Q_j$ are disjoint. Let
$C=\max\{2,\chi(H),h+m+1\}$. At least one of the $C$ labels
avoids all outside colour counts. Map the old colour class of $w$
to that label and inject the other proper classes into remaining
labels. If $\delta(H)\ge T(C)$, Theorem~\ref{thm:chromatic} fills
$H$ and validates the boundary without changing any given edge.

In particular, a vertex-minimal counterexample cannot have a bridge
with one side $H$ of minimum internal degree at least
$T(\max\{\chi(H),3\})$. Replace $H$ by a leaf and use minimality,
then the preceding absorption with $h=1,m\le1$. If the smaller
graph is $K_2$, set the bridge mixed directly; its outside endpoint
has counts $(1,1)$ and no other outside edge to check. This treats
the exceptional smaller graph without assuming it colourable.

\subsection{Several CRT regions and low-side ports}\label{sec:crtregions}

Let disjoint induced regions $H_i=G[W_i]$ cover $W$ and put
$U=V(G)\setminus W$. Fix all edges other than those internal to
a single region, including edges between different regions. Assume
all $U$--$U$ edges are legal with their actual fixed endpoint counts.
Given a proper $c$-colouring $j$ of the \emph{entire} graph $G[W]$,
let $h_v$ be the number of edges from $v\in W$ to $U$, and $m_v$
the number of mixed ones. For every nonempty colour class define
\[
 L_a=1+\max_{j(v)=a}(h_v+m_v),\qquad C=\sum_a L_a,
\]
using $c\ge2$; unused classes may be omitted.

\begin{theorem}[D2: CRT region absorption]\label{thm:crtregions}
If every $\delta(H_i)\ge T(C)$, the fixed colouring extends over
all regions. If each high vertex has at most $t$ neighbours in $U$,
the uniform sufficient condition is $\delta(H_i)\ge T(c(2t+1))$.
\end{theorem}
\begin{proof}
Reserve disjoint blocks of $L_a$ labels for the proper colour classes.
At $v$, outside counts prohibit at most $h_v+m_v<L_{j(v)}$ labels.
Choose an allowed label in its reserved block, independently at each
vertex. All high--high edges, including inter-region edges, still
have different labels. Feed all fixed incident contributions as the
external counts into Theorem~\ref{thm:chromatic} in each region.
Internal and inter-region edges distinguish both residues; high--low
edges avoid their fixed prohibited residues. Low--low edges and
all their endpoint counts persist. The estimate $C\le c(2t+1)$
gives the uniform bound, using additional unused labels if necessary.
\end{proof}

There is a convenient way to supply the low precolouring. For every
component $D$ of $G[U]$, attach a private placeholder leaf for each
actual high--low edge, at its low endpoint. Different ports use
different leaves even if their actual high endpoints agree. A legal
colouring of this placeholder graph transfers the placeholder edge
weights to the actual boundary edges, preserving each low vertex
degree in both colours. If $D$ is bipartite or a cactus, its
placeholder graph remains in a class covered by~\cite{GW23,GW24}.
The only $K_2$ placeholder case is a single low vertex with one
port; make that port mixed, with low counts $(1,1)$ and no low--low
edge. Thus no $K_2$ colouring is invoked. Inter-region edges can be
fixed arbitrarily before absorption. This gives complete infinite
families with global minimum degree one or two, rather than only
globally high-degree graphs.

\subsection{Trees attached to an ordinarily irregular core}\label{sec:irregularcore}

\begin{proposition}[D2: matched tree attachment]\label{prop:irregularcore}
Suppose $H$ is connected, has an edge, and its ordinary degrees
differ at the ends of every edge. Let $S\subseteq V(H)$ have a
matching to $V(H)\setminus S$ covering $S$. At each vertex of $S$
attach one disjoint arbitrary rooted tree by a new bridge. The
doubled expanded graph has a legal two-colouring.
\end{proposition}
\begin{proof}
Give matching edges and the new bridges weight 1, all other core
edges zero. Core counts are
\[
 (R,B)=\begin{cases}
 (2,2d_H(v)),&v\in S,\\
 (1,2d_H(v)-1),&v\text{ is a matching helper},\\
 (0,2d_H(v)),&\text{otherwise}.
 \end{cases}
\]
Matching red edges join degrees two and one. Blue degrees at helpers
are odd, elsewhere even; within either parity group ordinary degree
irregularity makes them unequal. Each tree bridge meets final core
counts $(2,2d_H(v))$. A singleton tree has counts $(1,1)$; a larger
tree has three mixed-port candidates to avoid the two bridge
exclusions by Lemma~\ref{lem:treeports}. Fill every tree independently.
\end{proof}

A single specified root always has such a matching, by choosing
one neighbour. Arbitrary many roots need not, and the resulting
expanded graph need not remain ordinarily locally irregular, so
unconditional repeated attachment does not follow. Moreover a
\emph{given} old colouring can be rigid. Take edges $vu,vw$, two
two-edge $u$--$z$ paths and three two-edge $w$--$z$ paths. Its
degrees at $v,u,w,z$ are $2,3,4,5$ and the five interiors have degree
two. All-mixed edges give a legal old colouring. A leaf at $v$
has no possible new weight preserving that colouring: weights zero,
one and two cause respectively a blue equality with $w$, both
equalities with $u$, and a red equality with $w$. Changing old
$vw$ to zero and giving the new edge weight 1 repairs the graph.
